\documentclass[11pt]{article}
\usepackage[T1]{fontenc}
\usepackage{lmodern}
\usepackage[letterpaper,margin=1in]{geometry}
\usepackage{amsmath,amssymb,booktabs,array,longtable,microtype}
\usepackage{xurl,listings}
\usepackage[hidelinks]{hyperref}
\lstset{basicstyle=\ttfamily\small,breaklines=true,columns=fullflexible,keepspaces=true}
\setlength{\parindent}{0pt}
\setlength{\parskip}{.35em}
\setlength{\emergencystretch}{2em}
\newcommand{\ind}{\mathbf 1}
\newcommand{\ReleaseStatus}{Version 2.1 supplement:
\url{https://github.com/rgillon/north-star-paper-a-supplement/tree/main/paper-a-v2.1}.
The repository index links the revised documents and complete archive.}
\hypersetup{pdftitle={North Star Paper A Version 2.1: Response to Referee Comments},
pdfauthor={Robert Gillon}}
\begin{document}
\begin{center}
{\Large\bfseries Response to Referee Comments}\\[.5em]
{\large Reconstructing Unions of Subgroups\\
from Cyclic Autocorrelation}\\[.4em]
Robert Gillon\quad$\cdot$\quad Version 2.1\quad$\cdot$\quad 22 November 2026
\end{center}

\section*{Scope of this response}
This response addresses the supplied Grok report and the subsequent
verification-report summary headed \emph{HOLDS WITH FIXES}.
The full attachment mentioned there was not among the supplied files;
this response concerns the quoted reports.

The comments prompted a fresh reading of the exact conductor argument,
an independent exact check of both cited cancellation examples, and
reproduction of the disputed supplement entries from a fresh extraction
of the delivered review archive. The revisions clarify a consequential
coordinate convention and correct genuine wording and invocation defects.
They do not strengthen the hypotheses of Theorem 10.4.
The set-level classification, exhaustive modulus criterion, and
construction results retain their separately stated scopes.

\section{The nontrivial-character condition in Theorem 10.4}
The reported concern is that Lemma 10.3 might supply only a nonempty
common zero set, permitting the all-\(u_i\) evaluations
\[
-2-4+6=0,\qquad
\frac14-\frac16-\frac1{12}=0.
\]
These identities are valid. The distinction needed in the proof is
between a \emph{nonempty common zero set} and a \emph{common zero at a
nontrivial character}. Lemma 10.3 states and proves the latter.

Let \(K=q_1q_2q_3\), put \(u_i=q_i-1\), and write a residual character as
\(\chi=\chi_1\otimes\chi_2\otimes\chi_3\). Its conductor \(c\mid K\)
means its character order. The local nonzero-shell sum is
\[
x_i(c)=\sum_{z\in C_{q_i}\setminus\{0\}}\chi_i(z)
=\begin{cases}
u_i,&q_i\nmid c,\\
-1,&q_i\mid c.
\end{cases}
\]
A trivial local character contributes \(q_i-1\). For a nontrivial
local character, its sum over the whole group is zero and its value at
the origin is one, leaving \(-1\) on the nonzero shell. Hence
\[
c>1\quad\Longleftrightarrow\quad
\text{some }\chi_i\text{ is nontrivial}
\quad\Longleftrightarrow\quad
\text{some }x_i=-1.
\]
In the kernel-exponent convention of equation (10.1), \(h_i\in\{0,1\}\)
and \(c=\prod_iq_i^{1-h_i}\). The trivial character has
\(h=(1,1,1)\); a nontrivial character has some \(h_i=0\).
Thus a \emph{zero coordinate} here means a zero \emph{kernel exponent},
not a zero shell sum. The all-\(u_i\) evaluation is precisely the
trivial character.

The revision inserts this dictionary and its character-sum proof at the
invocation of Lemma 10.3. It supplies an explicit justification of the
implication used there without an additional existence assertion.

\subsection*{Why Lemma 10.3 excludes the trivial-only case}
Write \(Z\) for the common zero conductors of the least-size sections
\(d_t\), and \(j\) for their first change on the primary axis.
The proof shows that the Fourier transform of \(T_j(g)\) vanishes
modulo \(p^{j+1}\) outside \(Z\). Inversion is legitimate because
\(p\nmid 2K\). If \(Z=\varnothing\), the function is zero modulo
\(p^{j+1}\); if \(Z=\{1\}\), it is constant modulo \(p^{j+1}\).

At a changing shell the possible values are
\[
p^j,\quad -p^j,\quad p^j-2p^b\ (0\le b<j),
\]
all nonzero modulo \(p^{j+1}\). This excludes \(Z=\varnothing\).
Under \(Z=\{1\}\), every \(d_t\) has zero population. If \(d_0\ne0\),
its two signs give value zero at the top residual shell, contradicting
constancy. Otherwise \(d_0=\cdots=d_j=0\), and the residual-origin
values are zero or \(-2p^c\), \(0\le c\le j\).

Comparison with the changing-shell values leaves only \(p=3\),
\(c=j\), and the positive residue \(p^j\). In that case the origin
fibers first reach size \(L\) at \(j+1\), bounding every later fiber
by \(L\). An earlier larger dispute would have zero partial sums
through \(j\), and sum zero or one at \(j+1\); it would give residue
zero or \(-p^j\), a contradiction. Thus only one disputed size remains.
Restricting to the dyadic subgroup of order \(L\) preserves homometry
and all \(d_t\). Reapplying the modular argument to the restricted sums
gives a changing-shell value \(p^j\) and a value zero at the residual
origin or top. This excludes the final trivial-only possibility.

Consequently \(Z\) contains a conductor \(c>1\). The remainder of
Theorem 10.4 uses the printed rank-one/rank-two exclusions and
exact-size population balance. No finite enumeration enters that proof.

\subsection*{Exact checks of the two cited examples}
Let \(T=\{i:q_i\mid c\}\). The table gives the complete spectra of
\[
P_1=-x_1-x_2+x_3
\quad\text{at }(q_1,q_2,q_3)=(3,5,7),
\]
and
\[
P_2=x_2x_3-x_1x_3-x_1x_2
\quad\text{at }(q_1,q_2,q_3)=(5,7,13).
\]
Dividing \(P_2\) by \(x_1x_2x_3\) gives the cited reciprocal expression.
\begin{center}
\begin{tabular}{@{}crrrr@{}}\toprule
\(T\) & \(c\), first example & \(P_1\) & \(c\), second example & \(P_2\)\\\midrule
\(\varnothing\)&1&0&1&0\\
\(\{1\}\)&3&3&5&90\\
\(\{2\}\)&5&5&7&-56\\
\(\{3\}\)&7&-7&13&-26\\
\(\{1,2\}\)&15&8&35&-1\\
\(\{1,3\}\)&21&-4&65&-1\\
\(\{2,3\}\)&35&-2&91&9\\
\(\{1,2,3\}\)&105&1&455&-1\\\bottomrule
\end{tabular}
\end{center}
Both polynomials vanish only at the trivial conductor.
The separate exact script
\path{verification/check_primary_conductor_coordinates.py} checks the shell-factor
values against direct signed residue polynomials reduced modulo
cyclotomic polynomials. This is a finite control of the reported
examples; the general implication follows from the character calculation
and the proof of Lemma 10.3.

\section{Reconciliation of the computational and artifact comments}
The delivered Version 2 archive contained 99 entries; all 98 files
listed in its checksum manifest matched after fresh extraction.
The following findings concern that archive. Accessibility at a public
branch tip is a separate obligation.

\subsection*{H100 height 3360 and the H99 target counts}
The \emph{separate supplementary audit} exists as the combined program
\path{verification/h99_independent_audit.py}.
Its \path{H100_rank} result freshly reproduced 174 endpoint inequalities,
435 actual grid targets, precision 3353, and maximum height 3360.
Corollary 17.2 is revised to name this file and its output explicitly.

Two programs explain the H99 counts. The original
\path{experiments/ppi_99_prime_pair.py} completes
\(3\cdot85=255\) reconstructions. The separate changed-data program
completes \(5\cdot85=425\). For the supplied data, its five targets
comprise both interval endpoints and three strictly interior grid points.
Thus the earlier phrase \emph{425 interior grid targets} is inaccurate:
the revision says \emph{425 actual dyadic-grid targets} and identifies
the implementation. Correspondingly, the H100 changed-data audit's
435 targets comprise 174 endpoints and 261 strict interior points.

The changed-data program imports no construction-verifier module, but
uses an independently checked H100 donor report as exact data. A fresh
run using a freshly regenerated donor report reproduced its mathematical
fields. The H99 error ratio is displayed as
\(0.3387485931919912\); the certification compares the underlying
rational exactly with \(339/1000<10/27\).
The proof for every target remains the interval, denominator and
parity argument, not the finite sample.

\subsection*{Paths, commands and the cut-moment row}
The explicit two-argument invocation succeeded. The separate audit's
no-argument defaults were genuinely stale: they pointed into an earlier
research directory and to the wrong donor-report directory.
The revised defaults use the bundle's
\path{experiments/} and \path{evidence/h100_independent.json}.
The explicit invocation remains:
\begin{lstlisting}
python3 verification/h99_independent_audit.py \
  experiments evidence/h100_independent.json
\end{lstlisting}
Section 11's command is changed from \texttt{python} to \texttt{python3}.

The current cut-moment row is reproducible with
\begin{quote}\small\ttfamily
python3 verification/verify\_cut\_moments.py
\end{quote}
Its JSON output matched \path{evidence/cut_moment_results.json} exactly:
85 moduli, 1,567 ideals, 503 membership cuts and 2,248 transfer
comparisons. The script and evidence file are named beside the relevant
claim. These bounded checks do not prove the universal theorem;
Theorem 7.1 and Proposition 7.2 supply that proof.
The older count of 73 comparisons, associated with an unrecovered
source formulation, remains unverified and is not adopted.

\subsection*{Software assumptions and the symbolic replay}
Fixed construction checks use exact integers and rational arithmetic;
decimal values are explanatory displays. The symbolic grid certificate
uses bounded integer arrays, with the stated coefficient bound
\(36\cdot13^6<2^{31}\), and exact rational nullspaces and polynomial
factorization. The documented environment uses Python 3.12.14,
NumPy 2.3.5 and SymPy 1.14.0, with assertions enabled.

The later report records a complete successful 24-shape replay, every
residual empty and all six table rows matched. That reported run is kept
distinct from the archived preparation run. Section 11 remains
\textsc{Computed}, because exhaustive implementation correctness is a
dependency even after successful replay. Theorems 1.1 and 17.3 retain
complete construction proofs with their explicitly identified fixed-data
certificates.

\section{Expository revisions requested in the Grok report}
\textbf{Quantifiers.}
The introduction makes explicit that individual rigidity fixes \(A\)
and quantifies over its possible mates \(B\), while modulus rigidity
quantifies over every admissible \(A\) and \(B\) at that modulus.
The complete modulus procedure may enumerate exponentially many ideals.

\textbf{Construction roadmap.}
A short roadmap explains the successive roles: upper and lower donors
supply complementary forced boundary values; modular helpers control
denominators while preserving annihilation; damping blocks reduce
interpolation errors; spectator factors realize the masses by actual
binary nested sections; and rank-dependent precision reduces the
required dyadic height. This separates the construction's logical
architecture from its fixed numerical parameters.

\textbf{Dyadic case and examples.}
The marginal-defect proof explicitly notes that its middle term vanishes
when \(p=2\), so the argument covers the dyadic axis. Section 8 directs
readers to Proposition 4.2 for the translation/complementation
interpretation of the radial examples.

\textbf{Possible arithmetic invariants.}
The open-question discussion identifies quantities suggested by the
sign graph: prime support of unanchored conductors, forced patterns of
Fourier zeros, and the prime-adic divisibility tests controlling adjacent
edges. These are research directions, not an asserted classification
by invariants of \(N\) alone. Structural arithmetic classification,
minimal nonrigid moduli and exact dyadic thresholds remain open.

\section{Supplement availability and review status}
\ReleaseStatus

The revised package assembles the original archive and the separately
written Version 2 verifiers, with the named evidence files and a checksum
manifest. The inspected initial public commit contained the earlier
archive and README, but not the later verifier additions. This observation
does not establish deletion. Public retrieval of the complete Version 2.1
package is recorded separately, with a pinned commit.

The initially requested generator-lattice rank and factor-cut moments
source notes remain unrecovered and unverified. The paper certifies its
printed conductor-grid rank theorem and exhaustive coefficient-moment
criterion without importing a compressed-generator complexity claim from
those notes.
\end{document}
